\documentclass{article} % For LaTeX2e

\usepackage{nips14submit_e,times}
\usepackage[spanish]{babel}
\usepackage{biblatex}
\usepackage{hyperref}
\usepackage{url}
\usepackage{listingsutf8}
\usepackage[utf8]{inputenc}
\usepackage{float}
\usepackage[T1]{fontenc}
\usepackage{graphicx}
\usepackage{placeins}
\usepackage{amsmath}
\usepackage{subfigure}
\usepackage{blindtext}
\usepackage{enumerate}
\usepackage{mwe}
\addbibresource{references.bib}
\usepackage[justification=centering]{caption}

\title{Inferencia Estadística \\ Situación 2: Comparación de Estimadores de \texorpdfstring{$\mu$}{mu}}

\author{
Jorge Andrés Jaramillo Neme \\
Universidad Autónoma de Occidente\\
\texttt{jorge\_and\.jaramillo@uao.edu.co} \\
\and
  \textbf{Jhonatan Andrés Tapia} \\
  Universidad Autónoma de Occidente \\
  \texttt{correo2@uao.edu.co} \\
  \and
  \textbf{Luis Fernando Meza} \\
  Universidad Autónoma de Occidente \\
  \texttt{correo3@uao.edu.co} \\
}



\nipsfinalcopy % camera-ready

\begin{document}
\maketitle

\section{Planteamiento}
Sea $X$ una población con
\[
E(X)=\mu, \qquad \operatorname{Var}(X)=\sigma^2,
\]
y sea $x_1,\dots,x_{2n}$ una muestra aleatoria simple de tamaño $2n$.
Se definen dos estimadores de la media poblacional $\mu$:
\begin{equation}
\bar X_1 \;=\; \frac{1}{2n}\sum_{i=1}^{2n} x_i,
\qquad
\bar X_2 \;=\; \frac{1}{n}\sum_{i=1}^{n} x_i.
\end{equation}

A continuación se calculan, para cada estimador: \emph{esperanza}, \emph{varianza}, \emph{sesgo} y \emph{error cuadrático medio (ECM)}.

\section{Resultados previos}
Usaremos que para constantes $a,b$ y variable aleatoria $Y$:
\[
E[aY+b]=a\,E(Y)+b,\qquad \operatorname{Var}(aY)=a^2\operatorname{Var}(Y),
\]
y para $Y_i$ \emph{independientes},
\[
\operatorname{Var}\!\Big(\sum_i Y_i\Big)=\sum_i \operatorname{Var}(Y_i).
\]
Además, por definición,
\[
\operatorname{ECM}(\hat\theta)=E\!\big[(\hat\theta-\theta)^2\big]
= \operatorname{Var}(\hat\theta) + \big(\operatorname{Sesgo}(\hat\theta)\big)^2.
\]

\section{Estimador \texorpdfstring{$\bar X_1$}{X1}}
\subsection{Esperanza}
\begin{align*}
E(\bar X_1)
&= E\!\left[\frac{1}{2n}\sum_{i=1}^{2n} x_i\right]
= \frac{1}{2n}\sum_{i=1}^{2n} E(x_i)
= \frac{1}{2n}\,(2n)\mu
= \mu.
\end{align*}

\subsection{Varianza}
\begin{align*}
\operatorname{Var}(\bar X_1)
&= \left(\frac{1}{2n}\right)^2
\operatorname{Var}\!\left(\sum_{i=1}^{2n} x_i\right)
= \left(\frac{1}{2n}\right)^2 \sum_{i=1}^{2n} \operatorname{Var}(x_i) \\
&= \left(\frac{1}{2n}\right)^2 (2n)\sigma^2
= \frac{\sigma^2}{2n}.
\end{align*}

\subsection{Sesgo}
\[
\operatorname{Sesgo}(\bar X_1)=E(\bar X_1)-\mu=\mu-\mu=0.
\]

\subsection{ECM}
Como el sesgo es cero,
\[
\operatorname{ECM}(\bar X_1)=\operatorname{Var}(\bar X_1)=\frac{\sigma^2}{2n}.
\]

\section{Estimador \texorpdfstring{$\bar X_2$}{X2}}
\subsection{Esperanza}
\begin{align*}
E(\bar X_2)
&= E\!\left[\frac{1}{n}\sum_{i=1}^{n} x_i\right]
= \frac{1}{n}\sum_{i=1}^{n} E(x_i)
= \frac{1}{n}\,(n)\mu
= \mu.
\end{align*}

\subsection{Varianza}
\begin{align*}
\operatorname{Var}(\bar X_2)
&= \left(\frac{1}{n}\right)^2
\operatorname{Var}\!\left(\sum_{i=1}^{n} x_i\right)
= \left(\frac{1}{n}\right)^2 \sum_{i=1}^{n} \operatorname{Var}(x_i) \\
&= \left(\frac{1}{n}\right)^2 (n)\sigma^2
= \frac{\sigma^2}{n}.
\end{align*}

\subsection{Sesgo}
\[
\operatorname{Sesgo}(\bar X_2)=E(\bar X_2)-\mu=\mu-\mu=0.
\]

\subsection{ECM}
\[
\operatorname{ECM}(\bar X_2)=\operatorname{Var}(\bar X_2)=\frac{\sigma^2}{n}.
\]

\section{Comparación y conclusión}
Ambos estimadores son \textbf{insesgados}:
\[
E(\bar X_1)=E(\bar X_2)=\mu \quad\Rightarrow\quad \text{Sesgo}=0.
\]
Sin embargo,
\[
\operatorname{Var}(\bar X_1)=\frac{\sigma^2}{2n}
\;<\;
\operatorname{Var}(\bar X_2)=\frac{\sigma^2}{n},
\]
por lo que
\[
\operatorname{ECM}(\bar X_1)=\frac{\sigma^2}{2n}
\;<\;
\operatorname{ECM}(\bar X_2)=\frac{\sigma^2}{n}.
\]
Entonces, bajo el criterio de \emph{menor ECM} (o mínima varianza entre insesgados), \textbf{$\bar X_1$ es el mejor estimador de $\mu$} porque utiliza más información (las $2n$ observaciones) y, en consecuencia, es más preciso.

\end{document}
